\subsection{The final assembly on one Wilson construction}
\label{ex-assembly}

The final assembly must use one hierarchy throughout. An operator-product
expansion for one subsequential limit and a reconstruction for another would
not establish the desired result. We first state its operator-product and trace requirements, attach the
proved trace result to that hierarchy, and then assemble the normalized
Wilson family, all of its character reconstructions and its physical stress
on the same limiting data.

\subsubsection{The operator-product and trace requirements}
Let \(\boldsymbol A=(A_i)_{i\in I}\) and
\(\boldsymbol B=(B_j)_{j\in J}\) be finite lists of graded fields and
write \(w(\boldsymbol A)=\sum_{i\in I}\deg A_i\). For each grade
\(\ell\), let \((E_{\ell,\alpha})_{\alpha\in\mathcal C_\ell}\)
be a finite collection of fields. The coefficients
\(c_{\boldsymbol A}^{\ell,\alpha}(\boldsymbol x;o)\) do not depend
on the phase. The cluster geometry at scale \(r>0\) is
\[
 |x_i-o|\le Mr,\qquad |x_i-x_{i'}|\ge\eta r\quad(i\ne i'),
\]
where \(M\ge0\) and \(\eta>0\). Spectators satisfy
\(|y_j-o|\ge c\) and \(|y_j-y_{j'}|\ge c\) for \(j\ne j'\),
where \(c>0\). For every \(\epsilon>0\), each fixed coefficient obeys
\begin{equation}\label{ex-ope-coefficient}
 |c_{\boldsymbol A}^{\ell,\alpha}(\boldsymbol x;o)|
 \le r^{\ell-w(\boldsymbol A)-\epsilon}
\end{equation}
at sufficiently small \(r\), uniformly in the indicated cluster geometry.
For each fixed truncation degree \(D\), the expansion requires
\begin{align}
 \bigg|S^s_{\boldsymbol A,\boldsymbol B}(\boldsymbol x,\boldsymbol y)
 -\sum_{\ell=0}^D\sum_{\alpha\in\mathcal C_\ell}
 c_{\boldsymbol A}^{\ell,\alpha}(\boldsymbol x;o)
 S^s_{E_{\ell,\alpha},\boldsymbol B}(o,\boldsymbol y)\bigg|
 \le r^{D+1-w(\boldsymbol A)-\epsilon}.
 \label{ex-ope-remainder}
\end{align}
Here \(s\) can be any character or the mixed state. The small-radius
threshold for the remainder is uniform over these choices and over the
allowed positions, but may depend on the fields, \(D,M,\eta,c,\epsilon\).
There is no assertion of one threshold working for all truncation degrees.
The quantitative passage from the local construction to these coefficient
and remainder estimates is Theorem~\ref{thm:ope-quantitative}, applied to
one prescription and its one limiting field hierarchy. Its hypotheses,
including the matched fields and the phase estimates, must hold for the
hierarchy used here; the numerical purchases alone do not establish them.

The trace requirement concerns a centered grade-four field \(\Theta\).
Set
\[
 T(r)=S^{\chi_*}_{\Theta,\Theta}
            ((\boldsymbol0,-r/2),(\boldsymbol0,r/2)).
\]
Its vacuum one-point functional vanishes, and the required asymptotic is
\begin{equation}\label{ex-trace-limit}
 \lim_{r\downarrow0}r^8\log^2(r^{-1})T(r)
       =\frac{3\dim_{\mathbb R}\mathfrak g}{\pi^4}>0.
\end{equation}
Theorem~\ref{thm:profile-exclusion} shows that this asymptotic implies the
required exclusion of finite complex linear combinations of mass-shell
tracks, including empty products. The class of profiles and their exact
normalization are given in Subsection~\ref{sec:trace-profiles}.
Theorem~\ref{thm:trace-selected-endpoint} supplies this entire trace
requirement, including the zero vacuum one-point functional and the
precise profile exclusion, for the selected funded hierarchy with its
own phase reconstruction and normalized physical stress. Its $\Theta$
and $\chi_*$ are the ones just used to define $T$: orthogonal covariance
identifies the radial correlation with this time-axis pair, whose
separation is $r$. The hierarchy and extraction are fixed before the
arbitrary positive null sequence used in that theorem. Thus the trace
assertion introduces no further subsequence of the hierarchy or new
normalization of its grade-four field. This discharges the trace part
of item~3 below whenever the joint construction uses these canonical
data; the operator-product and other joint-construction clauses retain
their separate suppliers.

\subsubsection{The phase mixture on joint boundary tests}

The mixture identity for separated Euclidean correlations must also hold
for the joint tempered boundary distributions. It is not enough to verify
it on products of one-variable tests and then use pointwise convergence.
Here the uniform slab bounds supply the needed common domination.

\begin{proposition}[Boundary barycenter]\label{prop:assembly-boundary-barycenter}
Use one mixed table and its character tables from
Subsection~\ref{ex-phase-selection}, with their common field union and the
slab bounds of Proposition~\ref{ex-phase-slab-bounds}. Let
$W_n^{\rm mix}$ and $W_n^\chi$ be the joint boundary distributions
constructed from those tables by Theorem~\ref{thm:joint-boundary}.
For every fixed finite field tuple $\boldsymbol v$ and every joint
Schwartz test $f\in\mathscr S((\mathbb R^4)^n)$, the map
$\chi\mapsto W_n^\chi(\boldsymbol v)(f)$ is continuous and
\begin{equation}\label{eq:assembly-boundary-barycenter}
 W_n^{\rm mix}(\boldsymbol v)(f)
       =\int_X W_n^\chi(\boldsymbol v)(f)\,d\nu(\chi).
\end{equation}
More generally, character continuity and the mixture identity hold for
the tested continuation at each fixed weakly ordered height tuple.
The same constants bound all these distributions by one continuous
Schwartz seminorm for the fixed tuple. No exceptional set of characters
or further extraction is used.
\end{proposition}
\begin{proof}
At arity zero every table and boundary functional is the constant one,
so consider $n>0$. Fix the labels and the strictly ordered imaginary
heights. The phase slab operators have the same norm bounds as the mixed
operators. Applying the padding and matrix estimate of
Subsection~\ref{sec:analytic-boundary} therefore gives a single exponent
$m\in\mathbb N$ and constant $C$ such that the continued scalar
functions $F_\chi$, including the mixed one, satisfy the gap bound
\eqref{eq:gap-power}. The bound is independent of $\chi$ and of the real
positions. The same observation applies after any fixed finite list of
label operations used in the continuation.

We first verify the mixture identity in the strict interior. Fix real
times $t_i$, spatial positions $x_i$ and ordered heights $y_i$. Along
$z_i(c)=iy_i+ct_i$, the functions $F_\chi(z(c),x;\boldsymbol v)$
are holomorphic in an open horizontal strip about the real axis: choose
the strip so that the finitely many imaginary gaps remain positive.
On each compact substrip they are bounded uniformly in $\chi$ by the
gap estimate. For $c=i\tau$ close to zero they are the separated
Euclidean functions, which depend continuously on $\chi$ by their
character representation. It follows that they depend continuously on
$\chi$ throughout the strip. Here is the normal-family argument for
this last step. Given $\chi_j\to\chi$, Cauchy's formula on slightly
larger discs gives uniform bounds for the derivatives of the functions
on each compact disc. Arzel\`a--Ascoli and a diagonal subsequence over
a countable exhaustion give a locally uniform holomorphic subsequential
limit. On the imaginary interval it equals $F_\chi$, by the already
known continuity there, and the one-variable identity theorem makes it
equal to $F_\chi$ everywhere. Every convergent subsequence has this
limit; if convergence at a point failed, the violating subsequence
would have a further locally uniform subsequence, a contradiction.
Thus the full sequence converges. Metrizability of $X$ turns this
sequential assertion into continuity.

The integral over $X$ is holomorphic in $c$: the same compact bound
permits interchanging the probability integral with contour integrals,
and Cauchy's formula on a disc then proves holomorphy. At imaginary
$c$ close to zero, the Euclidean mixture identity identifies it with
the mixed continued function. The identity theorem and $c=1$ give
\[
 F_{\rm mix}(t+iy,x;\boldsymbol v)
       =\int_X F_\chi(t+iy,x;\boldsymbol v)\,d\nu(\chi).
\]
The integrands are measurable jointly in $\chi$ and the real positions:
they are continuous in each variable separately, and on the separable
position space the continuous-position functions are pointwise limits
of their values on successively finer countable grids. This supplies
Borel simple approximations in the position variable with measurable
coefficients in $\chi$. The fixed-height bound and $f\in L^1$ permit
Fubini, so the same identity holds for the tested functions
$\psi_y(f)$ on every strictly ordered height tuple.

Use the same $m,C$ and direction $r$ in the boundary formula
\eqref{eq:edge-taylor} for every character. For $0<s\le1$ its integral
remainder is bounded by
$C\|D^{m+1}f\|_{L^1}/m!$, uniformly in the weakly ordered $y$ and
$\chi$. Its finite sum is bounded by
$C\sum_{j=0}^m\|D^jf\|_{L^1}/j!$. Both are integrable with respect
to the probability $\nu$, and their sum is a continuous Schwartz
seminorm, as proved in Theorem~\ref{thm:joint-boundary}. Substitute the
strict-interior mixture in this formula and apply Fubini to its
absolutely dominated remainder. This proves the identity at every weakly ordered $y$; at $y=0$ it is
\eqref{eq:assembly-boundary-barycenter}. The same domination and the
strict-interior character continuity proved above give continuity in
$\chi$ of each finite term and of the remainder by dominated
convergence. This proves the asserted character continuity and common
seminorm bound directly on joint tests.
\end{proof}

\subsubsection{A precise conditional assembly}
An internal extraction order is a pair \((u,p)\) of strictly increasing
maps \(\mathbb N\to\mathbb N\). Refining by \((v,q)\) gives
\((u\circ v,p\circ q)\). The full mesh order is \(u(N)=N\).
The need for a full mesh limit at the public endpoint is stronger than
convergence just along \(u\circ v\).

\begin{lemma}[Identifying the full mesh limit]\label{ex-full-mesh}
Fix a test and its insertion labels. Suppose that at every suitable period
\(p\) the full sequence \(W_{N,p}\) converges to some \(\ell_p\).
Suppose also that an ordered extraction supplies, at every suitable period
\(p\), numbers \(w_p\) such that \(W_{u(j),p}\to w_p\), where \(u\) is
strictly increasing; that the period \(p(q(k))\) is suitable for all
sufficiently large \(k\); and that \(w_{p(q(k))}\to s\). Then
\(\ell_p=w_p\) at every suitable period, and
\(\ell_{p(q(k))}\to s\), that is, the same volume limit \(s\) holds with
the full mesh limits.
\end{lemma}
\begin{proof}
Strict increase implies \(u(j)\ge j\) by induction, so \(u(j)\to\infty\).
The full mesh convergence therefore implies
\(W_{u(j),p}\to\ell_p\). To see uniqueness explicitly, for any
\(\epsilon>0\), take \(j\) sufficiently large that this term is within
\(\epsilon/2\) of each proposed limit. The triangle inequality gives
\(|\ell_p-w_p|<\epsilon\), hence equality. By hypothesis there is
\(k_0\) such that \(p(q(k))\) is suitable for \(k\ge k_0\); for those
\(k\) the equality just proved gives \(\ell_{p(q(k))}=w_{p(q(k))}\).
The two sequences of row limits therefore agree from \(k_0\) on and have
the same limit \(s\).
\end{proof}

This lemma explains how a full ultraviolet convergence result can be combined
with a hierarchy constructed by subsequence extraction. It does not obtain
full ultraviolet convergence from subsequential convergence: that is its
additional hypothesis.

\begin{assumption}[Joint construction input]\label{ex-joint-input}
For the fixed group \(G\) and each chosen positive blocking penalty
\(b\), there is one lattice family with the following properties.
Its insertions are continuous, gauge-invariant and uniformly local, and its
bare coefficients satisfy \eqref{ex-running-input}.
For every extraction order \((u,p)\), there is a refinement \((v,q)\)
and one set of limiting data as described in Section~\ref{ex-statement}
(graded fields with their injections, covering and token conditions, phase
data consisting of a compact metric space with a probability measure of
full support and a selected character, and phase and mixed correlations),
together with a selected relativistic theory as in
Definition~\ref{ex-reconstructed}, for which
all of the following hold simultaneously:
\begin{enumerate}
\item The mixed correlations satisfy \eqref{ex-row-target} along the
\emph{full} mesh sequence at every suitable period and
\eqref{ex-volume-target} along \(p\circ q\).
\item The phase correlations have the stated smoothness, symmetry,
continuity and mixture identity \eqref{ex-mixture}.
\item The fields and coefficients satisfy
\eqref{ex-ope-coefficient}--\eqref{ex-ope-remainder} and the centered
field satisfies \eqref{ex-trace-limit}.
\end{enumerate}
Every assertion refers to the same lattice family and the same selected
hierarchy on that refinement.
\end{assumption}

\begin{proposition}[Assembly from the joint input]\label{ex-conditional-existence}
Input~\ref{ex-joint-input} implies the full existence target stated in
Section~\ref{ex-statement}, including the operator-product and trace
properties above.
\end{proposition}
\begin{proof}
Fix \(G\) and choose \(b=1\). Select the family supplied by
Input~\ref{ex-joint-input} once. Its regularity is part of that input and
is unaffected by any subsequent choice of volume sequence. Now let
\(p:\mathbb N\to\mathbb N\) be an arbitrary strictly increasing
sequence. Apply the input to the extraction order \((\mathrm{id},p)\).
It yields \((v,q)\) and one hierarchy with all three simultaneous
properties. Take \(q\) as the public volume subsequence and retain this
hierarchy. Its ordered Wilson convergence uses the full mesh sequence by
item~1, so no assertion about a mesh subsequence has been substituted for
\eqref{ex-row-target}. Item~2 gives the phase and mixture conditions;
the selected theory is part of the same data. Item~3 supplies the remaining
two properties, with finite-profile exclusion following from
Theorem~\ref{thm:profile-exclusion}. The construction of the family preceded the choice of
\(p\), while \(q\) and the hierarchy were chosen after \(p\), exactly
as required.
\end{proof}

The proposition isolates the elementary logical projection to the public
target. The following assembly retains the additional normalization and
every-character assertions of the stronger headline. The all-period
supplier of Proposition~\ref{prop:fullmesh-selected-supplier} applies
to the retained family itself.

\subsubsection{Assembly of the construction-bound theorem}

\begin{theorem}[Assembly of the retained construction]
\label{thm:ex-construction-assembly}
The selected construction of Subsection~\ref{sec:public-run-binding} satisfies
Theorem~\ref{thm:ex-construction-headline}. In particular it
supplies Input~\ref{ex-joint-input}, and hence the public existence target.
The lattice family is fixed before the requested extraction order; all
limiting fields, phase decompositions, boundary distributions, local
products and the physical stress in the conclusion are constructed from
one limiting table on one refinement of that order.
\end{theorem}
\begin{proof}
Fix $G$ and the positive blocking penalty $b$. We describe the choices
in their order, then verify the four headline clauses.

\emph{The finite family.}
The local sections are chosen by Subsection~\ref{sec:local-sections}.
The common funded-history and observable purchases are made before
realizing the schedule: profiles and floors precede lengths; local
bounds and protection precede the final entry delays; the realized
clock and periods come last. The maxima in these choices are finite
at each stage, and subsequent enlargements preserve every previously
purchased lower bound. Subsection~\ref{sec:common-scale-selection}
therefore realizes all the requests with one schedule.
The annulus estimate and the retained path construction apply to its
one free horizon. The selected nonlinear return and the exact Wilson
startup then have their hypotheses on these paths and their original
quadratic data. Subsection~\ref{sec:surviving-interval-shooting} chooses
finite histories on every cutoff and period without changing that
schedule. The common-data identification of
Subsection~\ref{sec:public-run-binding} retains these histories together
with their transition witnesses and one family of fixed marked sections.

Proposition~\ref{prop:ex-c02-selected-record} places the complete
service record on that history, with its stored transitions and one
reference-normalized transfer. It is important to retain this transfer,
not merely the assertion that some valid transfer exists. The reference
normalization and the physical unit are exactly those of
Proposition~\ref{prop:public-normalized-observables}. Thus the public
family has the sections, the jet normalization and the exact inserted
density transport claimed in the headline. Its finite law is the
normalized Wilson law, by \eqref{eq:public-family-coordinates}; its
bare coefficient satisfies \eqref{eq:public-family-running}. All of
these choices have been made before an extraction order is mentioned.

\emph{One table and its field quotient.}
Now prescribe strictly increasing mesh and period maps $(u,p)$.
First retain the coupling and stationary limits along a refinement,
as in Subsection~\ref{sec:selected-history-limits}. Proposition~\ref{prop:fullmesh-selected-supplier} supplies the cap,
sweep and primary budgets for the actual retained insertions, with one
mesh threshold valid for every period. Thus the full-mesh bounds of
Theorem~\ref{thm:ex-fullmesh-fixedbounds-producer} and the common-history
Cauchy argument of Subsection~\ref{sec:fullmesh-correlation-cauchy}
apply to the same inserted densities. Theorem~\ref{ex-smooth-grid-limit}
gives their full mesh limits at every suitable fixed period. Separately,
Proposition~\ref{prop:ex-limits-extraction} and the finite-dimensional
extractions in Subsection~\ref{sec:finite-rank-kernel-transport} give a
common refinement for the table and all finite-grade projection data.
There are countably many grades, basis signatures and compact exhaustion
sets. The common diagonal extraction is made on these data; multilinearity
and the compact estimates then give the conclusions for every label and
test, rather than choosing a further subsequence for each test.
All previously obtained limits persist under this refinement.
Lemma~\ref{ex-full-mesh} identifies every extracted row with its full
mesh limit. Hence the mixed correlation formula has exactly the order
\eqref{eq:headline-wilson-limit}: full mesh at fixed suitable period,
then the chosen refinement of the prescribed period sequence.

The denominator, insertion and reflection identities are finite
identities of that same family. Proposition~\ref{prop:ex-limits-euclidean}
passes them to its table, with finite smeared positivity supplied by
Corollary~\ref{cor:ex-lattice-smeared-reflection}. Smoothness and the
local bounds pass by the smooth-limit and centered-moment estimates.
Take the null space of all separated mixed correlations with arbitrary
finite spectators, not only two-point correlations. The equation-kernel
relay and finite-rank construction identify precisely this null space.
The raw inclusion comparison is supplied by
Subsection~\ref{sec:actual-insertion-derivative-relay};
Proposition~\ref{prop:phase-field-inclusions} descends these comparisons,
and Proposition~\ref{prop:phase-field-union} forms the common union.
The calibrated physical unit gives Proposition~\ref{prop:ex-limits-unit}.
The derivative comparison and
Proposition~\ref{prop:phase-field-derivatives} give its compatible
positional derivatives. The conjugation and reflected adjoint are those
of Proposition~\ref{prop:phase-union-adjoint}. Each construction uses
the same all-spectator null spaces and the same table. This proves the
field and Wilson-limit assertions in item 1.

\emph{The actual character space and every character.}
Construct the mixed Euclidean Hilbert space from this table, then its
bounded slab operators and their compression to the heat-fixed space.
Subsection~\ref{ex-phases} proves that their algebra is commutative;
Proposition~\ref{ex-character-measure} gives its compact metrizable
character space and full-support vacuum probability. Its whole-word
character evaluations define the phase correlations, with the mixture
identity \eqref{ex-phase-barycenter}. Thus neither the phase space nor
its measure has been chosen independently of the Wilson limit.

Propositions~\ref{ex-phase-caps} and~\ref{ex-phase-slab-bounds}
give the common smoothness, caps and slab bounds in every character.
Lemma~\ref{ex-phase-original-null}, the finite detectors of
Proposition~\ref{ex-no-new-null} and their actual comparisons give the
same field null relations, including no additional null labels.
Proposition~\ref{prop:phase-inclusion-every-character} carries the same
inclusions into each phase. Reflection positivity, ordinary conjugation
and orthogonal covariance are those of
Sections~\ref{ex-phase-regularity} and~\ref{ex-reflection-transport}.
This verifies item 2 with its universal character quantifier.

For each character apply the complex heat construction, the joint
boundary theorem and Proposition~\ref{prop:heat-cups}. Their hypotheses
are the slab, differential and reflection properties just obtained,
which hold in every character. Propositions~\ref{prop:word-fields}
and~\ref{prop:word-motion} give the common invariant domain, cyclic
vacuum, adjoints, tempered fields and strongly continuous symmetry.
The calibrated label action descends by
Proposition~\ref{prop:lorentz-label-descent};
Theorem~\ref{thm:lorentz-covariance} and
Theorem~\ref{thm:spacelike-exchange} give covariance and locality.
Boundary detection and Proposition~\ref{prop:word-locality} retain
nonzero field labels. The positive-energy part of
Theorem~\ref{thm:field-spectrum} applies to every character.
The every-character vacuum result
Theorem~\ref{thm:ex-every-character-vacuum} supplies uniqueness of
the translation-invariant vacuum there; the almost-everywhere
selection theorem alone would not give this conclusion.
Finally Proposition~\ref{prop:assembly-boundary-barycenter} identifies
the phase mixture on all joint Schwartz tests with the boundary of the
mixed table. These are exactly the reconstructions in item 3.

\emph{One local product expansion and one physical stress.}
Apply the finite-menu and cluster estimates of
Subsection~\ref{sec:ope-history-service} to the stored transitions.
The terminal local projection and reference matching are performed on
the same insertion family and the same opened cells; their full
replicated difference, not products of separately transferred moments,
is estimated in Subsection~\ref{sec:reference-match}.
Lemma~\ref{lem:ope-selected-triangle} supplies the actual triangular
readouts, and Proposition~\ref{prop:ope-local-trial-comparison} supplies
the local-trial part of the physical-filtration argument. The
finite-rank projections above have exactly its all-test kernel.
Thus Proposition~\ref{prop:ope-match-assembly} supplies its matching
estimates on the one table. The grid argument and phase bounds give
Input~\ref{input:ope-raw-estimates}; the detector prescription is fixed
before choosing a phase, spectators or truncation degree. Its common
homogeneous version and Theorem~\ref{thm:ope-quantitative} give
\eqref{ex-ope-coefficient}--\eqref{ex-ope-remainder} simultaneously
for the mixed state and every character. Dependence of the small-radius
threshold on a fixed truncation degree is retained.

The finite-rank and comparison data used for the physical stress are
on this same quotient. Subsection~\ref{sec:actual-stress-normalization}
fixes its nonzero scalar normalization once across grades.
Corollary~\ref{cor:selected-ward-package} and
Proposition~\ref{prop:phase-ward-inheritance} give its conserved,
symmetric tensor and its translation and rotation charges in these
fields. Retain the distinguished character $\chi_*$ used in the
selected trace endpoint. Theorem~\ref{thm:trace-selected-endpoint}
applies to precisely this stress, hierarchy and character, with the
reconstruction just made. It centers the trace by this character's
one-point functional, proves the stated positive short-distance
coefficient and the finite-profile exclusion, and introduces no new
hierarchy or field normalization. This is item 4.

All four clauses consequently hold on the same refinement. Their
extraction-dependent choices were made only after the public family
was fixed. Forgetting the extra normalization and all-character
information supplies Input~\ref{ex-joint-input}; applying
Proposition~\ref{ex-conditional-existence} proves the weaker public
target without changing the family or the order of its limits.
\end{proof}
